\documentclass{article}
\usepackage{graphicx}
\usepackage{amsmath}
\usepackage{amsfonts}
\usepackage{bookmark}
\usepackage{booktabs}
\usepackage{exercises}
\graphicspath{ {} }
\usepackage{indentfirst}
\usepackage{amsthm}
\usepackage{yhmath}
\usepackage{geometry}
\usepackage{placeins}
\usepackage{csquotes}
\usepackage{pgfplots}
\newcommand*{\oldneg}{\mathord{\sim}}
\newcommand\showdiv[1]{\overline{\smash{)}#1}}
\newcommand\Mydiv[1]{%
\smash{\raise.3ex\hbox{$\big)$}}\mkern-8mu
        \overline{\enspace\strut#1}}
\setlength{\tabcolsep}{5pt}
\theoremstyle{definition}
\newtheorem{definition}{Definition}

\setlength{\parindent}{.5in}
\setlength{\parskip}{12pt}
\renewcommand{\arraystretch}{2}
\geometry{margin=1in}
\newcommand{\?}{\stackrel{?}{=}}
\pgfplotsset{width=7cm,compat=1.18}\usepgfplotslibrary{statistics}

\title{Test (pg. 605)}
\begin{document}
\maketitle

\section*{1}
First, by the pythagorean theorem, the hypotenuse of this triangle is $\sqrt{4^2+9^2}=\sqrt{97}$. Thus
\[ \sin\theta=\frac{4}{\sqrt{97}}=\frac{4\sqrt{97}}{97} \]
\[ \csc\theta=\frac{\sqrt{97}}{4} \]
\[ \cos\theta=\frac{9}{\sqrt{97}}=\frac{9\sqrt{97}}{97} \]
\[ \sec\theta=\frac{\sqrt{97}}{9} \]
\[ \tan\theta=\frac{4}{9} \]
\[ \cot\theta=\frac{9}{4} \]

\section*{2}
First we find the second leg of the triangle:
\[ \sqrt{11^2-8^2}=\sqrt{57} \] and thus
\[ \sin\theta=\frac{\sqrt{57}}{11} \]
\[ \csc\theta=\frac{11}{\sqrt{57}}=\frac{11\sqrt{57}}{57} \]
\[ \cos\theta=\frac{8}{11} \]
\[ \sec\theta=\frac{11}{8} \]
\[ \tan\theta=\frac{\sqrt{57}}{8} \]
\[ \cot\theta=\frac{8}{\sqrt{57}}=\frac{8\sqrt{57}}{57} \]

\section*{3}
First we find the second leg of the triangle:
\[ \sqrt{7^2-(\sqrt{13})^2}=\sqrt{36}=6 \] and thus
\[ \sin\theta=\frac{\sqrt{13}}{7} \]
\[ \csc\theta=\frac{7}{\sqrt{13}}=\frac{7\sqrt{13}}{13} \]
\[ \cos\theta=\frac{6}{7} \]
\[ \sec\theta=\frac{7}{6} \]
\[ \tan\theta=\frac{\sqrt{13}}{6} \]
\[ \cot\theta=\frac{6}{\sqrt{13}}=\frac{6\sqrt{13}}{13} \]

\section*{4}
\[ 260^\circ=260\times\frac{\pi}{180}=\frac{13\pi}{9} \]

\section*{5}
\[ -50^\circ=-50\times\frac{\pi}{180}=-\frac{5\pi}{18} \]

\section*{6}
\[ \frac{4\pi}{5}=\frac{4\pi}{5}\times\frac{180}{\pi}=144^\circ \]

\section*{7}
\[ \frac{8\pi}{3}=\frac{8\pi}{3}\times\frac{180}{\pi}=480^\circ \]

\section*{8}
We know that the reference angle of $150^\circ$ is $180^\circ-150^\circ=30^\circ$. Thus, because we know that $150^\circ$ lies in Quadrant II,
\[ \tan 150^\circ=-\tan 30^\circ=-\frac{1}{\sqrt{3}}=-\frac{\sqrt{3}}{3} \]

\section*{9}
We know that $-480^\circ$ is coterminal with $-480^\circ+360^\circ\times 2=240^\circ$, which has a reference angle of 
$240^\circ-180^\circ=60^\circ$. Because $240^\circ$ lies in Quadrant III, we therefore know that
\[ \sec(-480^\circ)=\sec 240^\circ=-\sec 60^\circ=-2 \]

\section*{10}
We know that $-\frac{5\pi}{3}$ is coterminal with $-\frac{5\pi}{3}+2\pi=\frac{\pi}{3}$. Thus
\[ \sin(-\frac{5\pi}{3})=\sin\frac{\pi}{3}=\frac{\sqrt{3}}{2} \]

\section*{11}
We know that $\frac{11\pi}{6}$ has a reference angle of $2\pi-\frac{11\pi}{6}=\frac{\pi}{6}$. Thus, because $\frac{11\pi}{6}$ lies in Quadrant IV,
\[ \cos\frac{11\pi}{6}=\cos\frac{\pi}{6}=\frac{\sqrt{3}}{2} \]

\section*{12}
In the range of $0^\circ\le \theta\le 180^\circ$, or $0\le\theta\le\pi$,
\[ \theta=\cos^{-1}1=0^\circ \] or 0.

\section*{13}
In the range of $-90^\circ <\theta<90^\circ$, or $-\frac{\pi}{2}<\theta<\frac{\pi}{2}$,
\[ \theta=\tan^{-1}\sqrt{3}=60^\circ \] or $\frac{\pi}{3}$.

\section*{14}
In the range of $-90^\circ\le\theta\le 90^\circ$, or $-\frac{\pi}{2}\le\theta\le\frac{\pi}{2}$,
\[ \theta=\sin^{-1}(-\frac{\sqrt{2}}{2})=\sin^{-1}(-\frac{1}{\sqrt{2}})=225^\circ \] or $\frac{5\pi}{4}$, as the reference angle of $225^\circ$ is
$225^\circ-180^\circ=45^\circ$, and because it lies in Quadrant III,
\[ \sin 225^\circ=-\sin 45^\circ=-\frac{1}{\sqrt{2}} \]

\section*{15}
In the range of $0^\circ\le\theta\le 180^\circ$, or $0\le\theta\le\pi$,
\[ \theta=\cos^{-1}(-\frac{\sqrt{3}}{2})=150^\circ \] or $\frac{5\pi}{6}$, as the reference angle of $150^\circ$ is
$180^\circ-150^\circ=30^\circ$, and because it lies in Quadrant II,
\[ \cos 150^\circ=-\cos 30^\circ=-\frac{\sqrt{3}}{2} \]

\section*{16}
First, we know that $B=180^\circ -(47^\circ+32^\circ)=101^\circ$. By the law of sines, we know that
\[ \frac{a}{\sin A}=\frac{c}{\sin C} \]
\[ \frac{a}{\sin 47^\circ}=\frac{12}{\sin 32^\circ} \]
\[ a=\frac{12\sin 47^\circ}{\sin 32^\circ} \]
\[ a\approx 16.561 \]
Again by the law of sines,
\[ \frac{b}{\sin B}=\frac{c}{\sin C} \]
\[ \frac{b}{\sin 101^\circ}=\frac{12}{\sin 32^\circ} \]
\[ b=\frac{12\sin 101^\circ}{\sin 32^\circ} \]
\[ b\approx 22.229 \]
Thus $a\approx 16.561,B=101^\circ,$ and $b\approx 22.229$.

\section*{17}
We can use the law of cosines:
\[ \cos A=\frac{b^2+c^2-a^2}{2bc} \]
\[ \cos A=\frac{12^2+17^2-24^2}{2(12)(17)} \]
\[ A\approx 110.517^\circ \]
and
\[ \cos B=\frac{a^2+c^2-b^2}{2ac} \]
\[ \cos B=\frac{24^2+17^2-12^2}{2(24)(17)} \]
\[ B\approx 27.923^\circ \]
and
\[ \cos C=\frac{a^2+b^2-c^2}{2ab} \]
\[ \cos C=\frac{24^2+12^2-17^2}{2(24)(12)} \]
\[ C\approx 41.560^\circ \]
Thus $A\approx 110.517^\circ,B\approx 27.923^\circ,$ and $C\approx 41.560^\circ$.

\section*{18}
This is the $SSA$ case with an acute angle. We know that $h=11\sin 63^\circ\approx 9.8$, and because $b<9.8$, these measurements do not represent a triangle.

\section*{19}
This is the $SAS$ case. We can use the law of cosines to solve it:
\[ c=\sqrt{a^2+b^2-2ab\cos C} \]
\[ c=\sqrt{23^2+19^2-2(23)(19)\cos 101^\circ} \]
\[ c\approx 32.508 \]
and thus, by the law of sines
\[ \frac{\sin A}{a}=\frac{\sin C}{c} \]
\[ \frac{\sin A}{23}=\frac{\sin 101^\circ}{32.508} \]
\[ \sin A=\frac{23\sin 101^\circ}{32.508} \]
\[ A\approx 43.989^\circ \] meaning $B\approx 180^\circ-(43.989^\circ+101^\circ)=35.011^\circ$. Thus
$A\approx 43.989^\circ,B\approx 35.011^\circ,$ and $c\approx 32.508$.

\section*{20}
This is the $SSS$ case, meaning we can use the law of cosines:
\[ \cos A=\frac{b^2+c^2-a^2}{2bc} \]
\[ \cos A=\frac{30^2+21^2-24^2}{2(30)(21)} \]
\[ A\approx 52.617^\circ \]
and
\[ \cos B=\frac{a^2+c^2-b^2}{2ac} \]
\[ \cos B=\frac{24^2+21^2-30^2}{2(24)(21)} \]
\[ B\approx 83.335^\circ \]
meaning $C\approx 180^\circ-(52.617^\circ+83.335^\circ)=44.048^\circ$. Thus
$A\approx 52.617^\circ,B\approx 83.335^\circ,$ and $C\approx 44.048^\circ$.

\section*{21}
We know that $C=180-(26+77)=77^\circ$. Thus, by the law of sines,
\[ \frac{a}{\sin A}=\frac{c}{\sin C} \]
\[ \frac{a}{\sin 26^\circ}=\frac{50}{\sin 77^\circ} \]
\[ a=\frac{50\sin 26^\circ}{\sin 77^\circ} \]
\[ a\approx 22.495 \]
and because $C=B$, $c=b$, so that $b=50$. Thus
$a\approx 22.495,b=50,$ and $C=77^\circ$.

\section*{22}
We know that the area is
\[ \frac{1}{2}bc\sin A \]
\[ =\frac{1}{2}\times 16\times 18\sin 81^\circ \]
\[ \approx 142.227 \] units squared.

\section*{23}
We know that the semi-perimeter is $\frac{1}{2}(a+b+c)=\frac{1}{2}(8+6+7)=10.5$, meaning the area is
\[ \sqrt{10.5(10.5-8)(10.5-6)(10.5-7)}\approx 20.333 \] units squared.

\section*{24}
We know that the semi-perimeter is $\frac{1}{2}(a+b+c)=\frac{1}{2}(25+24+38)=43.5$, meaning the area is
\[ \sqrt{43.5(43.5-25)(43.5-24)(43.5-38)}\approx 293.785 \] units squared.

\section*{25}
We know that the area is
\[ \frac{1}{2}ab\sin C \]
\[ =\frac{1}{2}\times 7\times 13\sin 111^\circ \]
\[ \approx 42.478 \] units squared.

\section*{26}
We know that the semi-perimeter is $\frac{1}{2}(a+b+c)=\frac{1}{2}(16+33+24)=36.5$, meaning the area is
\[ \sqrt{36.5(36.5-33)(36.5-16)(36.5-24)}\approx 180.931 \] units squared.

\section*{27}
We know that the area is
\[ \frac{1}{2}ac\sin B \]
\[ =\frac{1}{2}\times 12\times 18\sin 61^\circ \]
\[ \approx 94.459 \] units squared.

\section*{28}
We know that
\[ \tan 39^\circ=\frac{w}{80} \]
\[ w=80\tan 39^\circ \]
\[ w\approx 64.783 \] meters.

\section*{29}
I do not think the problem gave enough information to solve it. It does not say how far away the crane is from the building. The angle which the arm needs to be set will change
depending upon how far the crane is from the building.

\section*{30}
We can find this distance $d$ by using the law of cosines (noting that the angle opposite $d$ is $180^\circ-20^\circ=160^\circ$):
\[ d=\sqrt{40^2+25^2-2(40)(25)\cos 160^\circ} \]
\[ d\approx 64.065 \] miles.

\end{document}
